\documentclass[10pt,a4paper]{amsart}
\usepackage[margin=19mm]{geometry}
\usepackage{amsmath,amssymb,mathtools}
\usepackage[scr=rsfs,cal=euler]{mathalfa}
\usepackage{xfrac}
\usepackage[unicode,hidelinks,bookmarks=false]{hyperref}
\newcommand{\ie}{i.e.\ }
\newcommand{\cf}{cf.\ }
\newcommand{\resp}{respectively\ }
\newcommand{\Subsection}{\subsection}
\providecommand{\nobreakdash}{\nobreak\hbox{-}}
\setlength{\emergencystretch}{2em}
\multlinegap0pt
\usepackage{bm}
\usepackage{bbm}
\usepackage{dsfont}
\usepackage{xspace}
\usepackage{cancel}
\usepackage{mathtools}
\usepackage{amsthm}
\makeatletter
\@ifundefined{theorem}{\newtheorem{theorem}{Theorem}}{}
\@ifundefined{definition}{\newtheorem{definition}[theorem]{Definition}}{}
\@ifundefined{lemma}{\newtheorem{lemma}[theorem]{Lemma}}{}
\makeatother
\def\Eh{{\hat{\mathbb{E}}}}
\def\Vh{{\hat{\mathbb{V}}}}
\title[Strong law of large numbers for $\varphi$-sub-Gaussian rando]{Strong law of large numbers for $\varphi$-sub-Gaussian random variables under sub-linear expectation spaces}
\author{Nyanga Honda Masasila, Istv\'an Fazekas}
\date{Source: arXiv:2602.18175v1 (CC BY 4.0)}
\begin{document}
\maketitle

\noindent\emph{Source and licence.} Strong law of large numbers for $\varphi$-sub-Gaussian random variables under sub-linear expectation spaces. Version arXiv:2602.18175v1 (CC BY 4.0), distributed under the Creative Commons Attribution 4.0 International licence, as recorded in the arXiv licence metadata for this exact version and shown on the abstract page https://arxiv.org/abs/2602.18175v1. Full rights notice: licenses/arxiv-2602.18175-CC-BY-4.0.txt.\par\smallskip

\noindent\textit{Editorial scope.} Counted unit: the Theorem whose source label is \texttt{MainTheorem} in the article (source0.tex lines 391--417, source label \texttt{MainTheorem}), delivered together with the article's own complete original proof of it (source0.tex lines 418--469, one \texttt{proof} environment of 52 lines). No mathematics is written, repaired or re-derived: statement and proof are the article's own text line for line. Required context retained uncounted: upGauss (lines 295--307); lowGauss (lines 303--305); mainLemma (lines 326--334); tau (lines 370--382). Every definition, display and label of those lines is retained unchanged. Omitted from this package and not redistributed: 1--294, 306--325, 335--369, 383--390, 470--643. Nothing inside a retained excerpt is omitted or altered: every retained body is delivered exactly as the article's lines are written, so no omission receipt of any kind accompanies this package. Citations: the retained excerpts carry no citation command at all, so no bibliography entry of the article is transcribed and no reference is redistributed. Reference adaptation: none was needed and none was made; every reference inside the delivered unit resolves inside it, so no unresolved reference or citation is produced. Printed slips kept literal: no printed word, symbol, hypothesis or number of the counted statement or of its proof was altered. Layout and class: the article's own class is \texttt{article} with packages including openwork, amsmath, amssymb, bm, amsthm, amscd, amsfonts, graphicx, color; the wrapper is the lane's pinned \texttt{amsart} editorial wrapper at 10pt on A4, which restates the theorem-like environments the retained text needs and adds the article's own macro definitions that the retained text needs (\texttt{\textbackslash Eh}, \texttt{\textbackslash Vh}), with extra wrapper packages (bm, bbm, dsfont, xspace, cancel, mathtools). The delivered numbering is the unit's own. Assets: no retained span contains a figure environment, a graphics-inclusion command, a PDF-page inclusion, a table or a TikZ picture; no image file is copied into this package, the package declares no asset dependency and no image file is redistributed with it. Licence: version arXiv:2602.18175v1 of this article is distributed under the Creative Commons Attribution 4.0 International licence, as recorded in the arXiv licence metadata for this exact version and shown on the abstract page https://arxiv.org/abs/2602.18175v1. A full rights notice is delivered at licenses/arxiv-2602.18175-CC-BY-4.0.txt. Characters: every retained excerpt is pure ASCII, so the delivered engine, XeLaTeX with its default Latin Modern text font, reports no missing character.
\begin{definition}
Let $\varphi$ be a quadratic $N$-function.
A random variable $\xi$ is said to be $\varphi$-sub-Gaussian under the sub-linear expectation $\Eh$ if there exist
constants  $a > 0$, $\overline{m}$, and $\underline{m}$, such that
\begin{equation}   \label{upGauss}
\Eh e^{\lambda(\xi - \overline{m}}) \le e^{\varphi(a\lambda)}, \quad {\text {for}} \quad \lambda>0,
\end{equation}
and 
\begin{equation}   \label{lowGauss}
\Eh e^{\lambda(\xi - \underline{m}}) \le e^{\varphi(a\lambda)}, \quad {\text {for}} \quad \lambda<0.
\end{equation}
$a > 0$, $\overline{m}$, and $\underline{m}$ are parameters of the $\varphi$-sub-Gaussian variable.
\end{definition}

\begin{lemma} \label{mainLemma}
Assume that \eqref{upGauss} and \eqref{lowGauss} are satisfied.
Then for $\varepsilon>0$ we have
\begin{equation}
\Vh\left( \{ \xi-\overline{m} > \varepsilon \} \cup \{ \xi-\underline{m} < -\varepsilon \} \right)  
\le 2 e^{-\varphi^*\left(\varepsilon/a\right)}.
\end{equation}

\end{lemma}

\begin{definition} \label{tau}
For a $\varphi$-sub-Gaussian random variable $\xi$ with fixed $\overline{m}$ and $\underline{m}$
 let $\tau_\varphi(\xi)$ be defined as
\begin{eqnarray*}
\tau_\varphi(\xi)
=
\inf\Big\{
a \ge 0 \, :\, 
\Eh e^{\lambda(\xi - \overline{m}}) \le e^{\varphi(a\lambda)}, \ \ {\text {for}} \ \ \lambda>0, \quad 
\text{and} \ \ \Eh e^{\lambda(\xi - \underline{m}}) \le e^{\varphi(a\lambda)}, \ \ {\text {for}} \ \ \lambda<0  
\Big\}.
\end{eqnarray*}
\end{definition}

\begin{theorem} \label{MainTheorem}
For some $p>1$, let $Z_n, n\ge1$, be a sequence of  $\varphi_p$-sub-Gaussian random variables with parameters
$\overline{m}$ and $\underline{m}$.
Let $\tau_\varphi(Z_n)$ be defined according to Definition \ref{tau}.
If there exist positive numbers $c$ and $\alpha$ such that for every natural number $n$, 
the condition $\tau_{\varphi_p}\!( Z_n )\le c\, n^{-\alpha}$ is satisfied, then
\begin{equation}
\label{upper}
\Vh\left(
\left\{\liminf_{n\to\infty}{Z_n} < \underline{m}\right\}
\cup
\left\{\limsup_{n\to\infty}{Z_n} > \overline{m}\right\}
\right)
= 0
\end{equation}
and 
\begin{equation}
\label{lower}
v\left(
\underline{m}
\le \liminf_{n\to\infty}{Z_n}
\le \limsup_{n\to\infty}{Z_n}
\le \overline{m}
\right)
= 1.
\end{equation}
\end{theorem}

\begin{proof}
Since $V(\cdot)$ and $v(\cdot)$ are conjugate to each other, \eqref{lower} is equivalent to \eqref{upper}.

By Lemma \ref{mainLemma}, for $\varepsilon>0$ we have
\begin{equation}  \label{V}
\Vh\left( \{ Z_n-\overline{m} > \varepsilon \} \cup \{ Z_n-\underline{m} < -\varepsilon \} \right)  
\le 2 \exp\left(-\varphi_q\left(\frac{\varepsilon}{\tau_{\varphi_p}(Z_n)}\right)\right) ,
\end{equation}
where $1/p+1/q=1$ and we applied that  $\varphi_q=\varphi_p^{*}$.
%
By the condition $\tau_{\varphi_p}\!( Z_n )\le c\, n^{-\alpha}$, 
for all sufficiently large $n$ we have $\frac{\varepsilon}{\tau_{\varphi_p}(Z_n)}>1$, so using the definition of
$\varphi_q$, we obtain that the right-hand side of inequality \eqref{V} is majorated by
$C_0\exp\!\left(-K n^{q\alpha}\right),$ 
where $C_0= 2 \exp\!\left(\frac{1}{q}-\frac{1}{2}\right)$,
and $K=\frac{1}{q}\left(\frac{\varepsilon}{c}\right)^q$.

We now show that the series
\[
\sum_{n=1}^{\infty} \Vh\left( \{ Z_n-\overline{m} > \varepsilon \} \cup \{ Z_n-\underline{m} < -\varepsilon \} \right) 
\]
is convergent. 
It suffices to prove the convergence of
\[
\sum_{n=1}^{\infty}\exp\!\left(-K n^{\beta}\right),
\qquad {\text{where}} \ \ \beta=q\alpha>0.
\]
The function $f(x)=\exp(-Kx^\beta)$ is positive and eventually decreasing.
By the integral test,
\[
\sum_{n=1}^{\infty}\exp(-K n^\beta)
\le
\int_{0}^{\infty}\exp(-Kx^\beta)\,dx.
\]
Performing the substitution $t=Kx^\beta$, we obtain
\[
\int_{0}^{\infty}\exp(-Kx^\beta)\,dx
=
\frac{1}{\beta}K^{-1/\beta}
\int_{0}^{\infty} t^{1/\beta-1}e^{-t}\,dt
=
\frac{1}{\beta}K^{-1/\beta}
\Gamma\!\left(\frac{1}{\beta}\right)
<\infty.
\]
Now, by the Borel-Cantelli lemma, it follows that the $\Vh$ measure is zero that infinitely many of the events
$$
\left( \{ Z_n-\overline{m} > \varepsilon \} \cup \{ Z_n-\underline{m} < -\varepsilon \} \right)
$$
occur.
As it is true for any $\varepsilon>0$, so \eqref{upper} is true.
\end{proof}

\end{document}
